% !TeX program = xelatex
% UTF-8. Compile in this directory with XeLaTeX.
% Define \SSInNotebook before \input to include in an existing XePersian document.
% Source: Google Drive Notebook/SS, handwritten pages 1--2.
% pic01.jpg = 20261018_200457752.jpg; 1XPNn99A7G1ZCiLqPpSHSCUVs0BvD_rvL
% pic02.jpg = 20261018_200503007.jpg; 1vn2jPXBaSdY8C98yVVp5YpasGVADcNYB
% Source wording, abbreviations, and findings are retained without clinical revision.
% Bold type represents highlighting; arrows preserve the source relationships.
\ifdefined\SSInNotebook
  \let\SSfinish\relax
\else
  \def\SSfinish{\end{document}}
  \documentclass[12pt,a4paper]{article}
  \usepackage[margin=20mm]{geometry}
  \usepackage{amsmath,amssymb}
  \usepackage{graphicx}
  \usepackage{enumitem}
  \usepackage{xepersian}
  \IfFontExistsTF{Amiri}{\settextfont{Amiri}}{\settextfont{DejaVu Sans}}
  \setlatintextfont{DejaVu Serif}
  \setlength{\parindent}{0pt}
  \setlength{\parskip}{5pt}
  \setlength{\emergencystretch}{3em}
  \setlist[itemize]{itemsep=3pt,topsep=3pt}
  \begin{document}
\fi
\begingroup
\newcommand{\SSarrow}{\lr{\(\longleftarrow\)}}
\newcommand{\SSrange}[2]{\lr{#1--#2}}
\newcommand{\SSheader}[1]{%
  \begin{center}
    {\Large\bfseries\lr{Serum Sickness} ـ ۸۲}\par
    \small\lr{Subject: Serum Sickness}\quad
    \lr{Year:}\quad\lr{Month:}\quad\lr{Date:}\par
    صفحهٔ #1
  \end{center}}

\SSheader{۱}

\subsection*{اتیولوژی}
واکنش ازدیاد حساسیت \textbf{تیپ \lr{III}}.

دارو یا آنتی‌سرم: به‌عنوان آنتی‌ژن شناخته شدن و علیه آن \lr{Ab}
\SSarrow{} تشکیل \textbf{کمپلکس \lr{Ag-Ab}}
\SSarrow{} رسوب در عروق و ارگان‌های تصفیه‌کننده
\SSarrow{} آبشار کمپلمان را فعال می‌کند
\SSarrow{} نوتروفیل‌ها فعال
\SSarrow{} افزایش نفوذپذیری عروق
\SSarrow{} آسیب بافتی.

\subsection*{اپیدمیولوژی}
امروزه کمتر دیده می‌شود.

عوامل شایع \lr{SS}:
\begin{itemize}
  \item فرآورده‌های خونی، آنتی‌ونوم‌ها، آنتی‌تیموسیت گلوبولین.
  \item برخی داروها \SSarrow{}
  \textbf{پنی‌سیلین، سولفونامیدها، تیازیدها}، \ldots
\end{itemize}

\subsection*{تظاهرات بالینی \lr{SS}}
\textbf{\SSrange{7}{21} روز} پس از ورود دارو یا عامل حساسیت‌زا ایجاد می‌شود.
در افرادی که از قبل حساس شده‌اند: \textbf{\SSrange{1}{4} روزه}.

\textbf{تظاهرات کلاسیک: تب، آرترالژی پلی‌آرتیکولار، لنفادنوپاتی، درگیری پوستی.}

\begin{itemize}
  \item انواع درگیری پوست: اورتیکاریا، آنژیوادم، اریتم مولتی‌فرم،
  راش موربیلی‌فرم، پالپبل پورپورا؛
  \lr{serpiginous rash} \SSarrow{} روی دست و پا.
  \item درگیری‌های کمتر شایع: کاردیت، گلومرولونفریت، سندرم گیلن‌باره،
  انسفالومیلیت، نوریت پریفرال.
\end{itemize}

\subsection*{یافته‌های آزمایشگاه}
\begin{itemize}
  \item \textbf{افزایش \lr{ESR}، ترومبوسایتوز}.
  \item \textbf{افت کمپلمان} \lr{(\(C_3, C_4\))}.
  \item بیوپسی پوست \SSarrow{} \textbf{تجمع آنتی‌بادی}:
  \lr{\(\mathrm{IgM},\ \mathrm{IgA},\ \mathrm{IgE}\ \text{or}\ C_3\)}.
  \item شواهد \SSarrow{} \textbf{هماچوری و پروتئین‌اوری}.
\end{itemize}

\medskip
{\footnotesize شماره‌های چاپی سطرهای برگه: \lr{1, 2, 3, 4, 5, 6, 7, 8}.\par}

\clearpage
\SSheader{۲}

\subsection*{درمان}
\begin{itemize}
  \item خودمحدودشونده / \textbf{ظرف \SSrange{1}{2} هفته}.
  \item \textbf{آنتی‌هیستامین} (برای خارش)،
  \textbf{\lr{NSAID}} (برای درد مفاصل) و (تب).
  \item در موارد شدید \SSarrow{} گلوکوکورتیکوئید سیستمیک.
\end{itemize}

\medskip
{\footnotesize شماره‌های چاپی سطرهای برگه: \lr{1, 2, 3, 4, 5, 6, 7, 8, 9}.\par}

% Original images retain all handwriting, printed labels, and layout marks.
\IfFileExists{images/pic01.jpg}
  {\def\SSimagepath{images/}}
  {\def\SSimagepath{pediatrics/ped-immunology/SS/images/}}
\newcommand{\SSoriginal}[1]{%
  \clearpage
  \begin{center}
    \textbf{تصویر اصلی \lr{#1}}\par\smallskip
    \includegraphics[width=\textwidth,height=0.91\textheight,keepaspectratio]{\SSimagepath pic#1.jpg}
  \end{center}}
\SSoriginal{01}
\SSoriginal{02}

\endgroup
\SSfinish
